%! Introduction to Experimental Design — Unit 3, Lesson 3.5 — Student Packet Cover Sheet
\documentclass[10pt]{article}
\usepackage{apstats-article}
\usepackage{apstats-boxes}
\usepackage{ltablex}
\keepXColumns

\begin{document}

% Full-bleed navy banner
\begin{tikzpicture}[remember picture, overlay]
  \fill[navy] (current page.north west)
    rectangle ([yshift=-0.9in]current page.north east);
\end{tikzpicture}
\vspace{-0.8in}

\begin{center}
  {\color{white}\textbf{\LARGE AP Statistics}}\\[4pt]
  {\color{skymid}\large\textbf{Unit 3: Collecting Data}}\\[6pt]
  {\color{white}\textbf{\large Lesson 3.5 \quad Introduction to Experimental Design}}
\end{center}

\vspace{0.22in}
\namedateperiod
\vspace{0.18in}

\begin{learningtargetbox}
\small
\begin{enumerate}[leftmargin=1.5em, itemsep=5pt]
  \item \ldots{} identify the experimental units, factor(s), treatments, response variable, and confounding variables in an experiment.
  \item \ldots{} describe the four principles of a well-designed experiment: comparison, random assignment, replication, and control.
  \item \ldots{} compare completely randomized, block, and matched-pairs designs.
\end{enumerate}
\end{learningtargetbox}

\vspace{0.14in}

\begin{tocbox}
\small
\renewcommand{\arraystretch}{1.5}
\begin{tabularx}{\linewidth}{@{} c l X r @{}}
\rowcolor{sky}
\textbf{\#} & \textbf{Component} & \textbf{Description} & \textbf{Score} \\
\midrule
1 & Warm-Up & Spiral review + experiment vs.\\ observational study & \blank{1.2cm} \\
2 & Guided Notes & Components of experiments; CRD, block, matched pairs & \blank{1.2cm} \\
3 & Group Activity & Identifying components and design type & \blank{1.2cm} \\
4 & Exit Ticket & Headache medication experiment & \blank{1.2cm} \\
5 & Homework & Four experiment descriptions; components and design & \blank{1.2cm} \\
\midrule
  & \multicolumn{2}{r}{\textbf{Total}} & \blank{1.2cm} \\
\end{tabularx}
\end{tocbox}

\end{document}
